% =========================================================================
% Seguimiento de riesgos — 26-07-2026 (S1, primera retrospectiva).
% Documento independiente de la memoria; los seguimientos nuevos se crean
% copiando TEMPLATE.tex (ver README.md de esta carpeta).
% Diseño heredado de la plantilla de seguimiento de Gestión de Proyectos
% del autor, adaptado a la configuración global de la memoria (preámbulo
% compartido). Los campos del formulario son los de la plantilla propia
% del TFG (TEMPLATE), no los del original: severidad en vez de categoría,
% nuevo impacto explícito y el estado del cap. 7 con sus cinco valores.
% Compilar DESDE ESTA CARPETA (2 pases, sin biber):
%   pdflatex --enable-installer -interaction=nonstopmode -file-line-error
%     -output-directory=build 2026-07-26-S1.tex
% =========================================================================
\documentclass[11pt,a4paper]{article}

% =========================================================================
% Preámbulo COMPARTIDO de los documentos independientes de riesgos
% (docs/thesis/plan-riesgos/ y docs/thesis/seguimiento-riesgos/).
%
% Estos documentos NO forman parte de main.tex (la memoria): son entregables
% sueltos en clase article, con el diseño heredado de los .tex de Gestión de
% Proyectos del autor, pero adaptados a la configuración global de la
% memoria (Carlito, márgenes 2,5 cm, parskip 6 pt, títulos de la guía).
% Un solo preámbulo para toda la familia: un ajuste de estilo aquí se
% propaga a todos los seguimientos futuros (uno por sprint hasta S37).
%
% RUTAS: los \input son relativos a la carpeta del documento, no a esta;
% compilar SIEMPRE desde la carpeta del .tex (plan-riesgos/ o
% seguimiento-riesgos/):
%   pdflatex --enable-installer -interaction=nonstopmode -file-line-error
%     -output-directory=build <documento>.tex        (2 pases; sin biber)
% =========================================================================

% --- Codificación (idéntico a packages.tex de la memoria) ---
\usepackage[utf8]{inputenc}
\usepackage[T1]{fontenc}

% --- Idioma: solo español (aquí no hay resumen trilingüe) ---
% es-noshorthands por coherencia con la memoria: acentos en UTF-8, sin
% caracteres activos (el conflicto con biblatex-apa no aplica aquí, pero
% mantener una única configuración de babel evita sorpresas al copiar texto
% entre la memoria y estos documentos).
\usepackage[spanish, es-tabla, es-noshorthands]{babel}

% --- Márgenes: 2,5 cm laterales e inferior, como la memoria ---
% El margen superior se reparte para alojar la cabecera enmarcada del
% diseño de referencia (2 cm de caja + separación) sin comerse el cuerpo:
% el texto arranca a 4 cm del borde y el ancho de texto sigue siendo 16 cm.
\usepackage[a4paper, left=2.5cm, right=2.5cm, bottom=2.5cm,
            top=1.2cm, includehead, headheight=2.2cm, headsep=0.6cm]{geometry}

% --- Tipografía: Carlito, la misma decisión que en style.tex ---
% (clon métrico libre de Calibri; pdflatex no puede usar la real).
% No se replica el arreglo de versalitas de style.tex: aquí no hay
% bibliografía APA y ningún texto usa \textsc.
\usepackage[sfdefault]{carlito}

% --- Espaciado: interlineado sencillo + 6 pt entre párrafos (guía) ---
\usepackage[skip=6pt]{parskip}

% --- Tablas ---
% longtable: el registro y el plan de respuesta superan una página.
% AVISO ESPERADO: "ignored error: Infinite glue shrinkage" en cada salto de
% página de una longtable (bug conocido e inofensivo, latex3/latex2e#1907;
% ver el comentario largo en packages.tex de la memoria).
\usepackage{array}
\usepackage{longtable}
\usepackage{tabularx}
% \RaggedRight (ragged2e) en vez de \raggedright en columnas muy estrechas:
% el \raggedright de LaTeX desactiva de facto la partición silábica y una
% palabra larga («Identificado») desborda una columna de 1,2 cm; la
% variante de ragged2e mantiene el guionado.
\usepackage{ragged2e}
% Las columnas X se componen como m{...}: el contenido de cada celda queda
% centrado verticalmente (aspecto de formulario del diseño de referencia).
\renewcommand{\tabularxcolumn}[1]{m{#1}}
% Tipos de columna del diseño de referencia:
%   L/C = ancho fijo alineado izquierda/centro; Y/Z = X alineada izq./centro.
% En columnas estrechas se prefiere Y (raggedright) a X justificada: la
% justificación en 3-4 cm abre ríos de blancos.
\newcolumntype{L}[1]{>{\raggedright\arraybackslash}m{#1}}
\newcolumntype{C}[1]{>{\centering\arraybackslash}m{#1}}
\newcolumntype{Y}{>{\raggedright\arraybackslash}X}
\newcolumntype{Z}{>{\centering\arraybackslash}X}
% Filas con aire en todas las tablas (el diseño de referencia lo fijaba
% tabla a tabla; aquí una sola vez).
\renewcommand{\arraystretch}{1.25}

% --- Utilidades ---
\usepackage{xcolor}    % \markfield (resaltado) y \todonote (macros.tex)
\usepackage{amssymb}   % \square, \blacksquare, \circlearrowright
\usepackage{xspace}    % lo necesita \bukit (macros.tex)
\usepackage{enumitem}
\setlist[itemize]{leftmargin=*, topsep=2pt, itemsep=2pt}
\setlist[enumerate]{leftmargin=*, topsep=2pt, itemsep=2pt}

% --- Títulos según la guía de estilo, como en style.tex de la memoria ---
% (sin \chapter: la clase article no lo tiene; el resto, idéntico).
\usepackage{titlesec}
\titleformat{\section}
  {\normalfont\bfseries\fontsize{12pt}{14.4pt}\selectfont\raggedright}
  {\thesection.}{0.5em}{\MakeUppercase}
\titleformat{\subsection}
  {\normalfont\bfseries\fontsize{12pt}{14.4pt}\selectfont}
  {\thesubsection.}{0.5em}{}
\titleformat{\subsubsection}
  {\normalfont\itshape\fontsize{12pt}{14.4pt}\selectfont}
  {\thesubsubsection.}{0.5em}{}

% --- Paginación dinámica (regla estricta del autor, 26-07-2026) ---
% PROHIBIDO el salto de página manual también en estos documentos: el
% espacio lo garantizan estos ganchos (o un \Needspace* local ante un bloque
% indivisible), nunca un \newpage. Los .tex de referencia sí usaban
% \newpage; es la única desviación deliberada respecto a ellos.
\usepackage{needspace}
\AddToHook{cmd/section/before}{\Needspace*{5\baselineskip}}
\AddToHook{cmd/subsection/before}{\Needspace*{5\baselineskip}}
\AddToHook{cmd/subsubsection/before}{\Needspace*{4\baselineskip}}
\widowpenalty=10000
\clubpenalty=10000
\displaywidowpenalty=10000
\raggedbottom
% Estos formularios componen mucho texto en columnas de 2-4 cm: se concede
% a TeX un margen de estirado de emergencia antes de desbordar la línea
% (evita Overfull \hbox menores sin tocar el contenido).
\setlength{\emergencystretch}{1em}

% --- Datos globales y macros de la memoria (única fuente: no duplicar) ---
% =========================================================================
% Variables globales del documento.
% Cambiar un dato AQUÍ lo actualiza en todo el documento (portada,
% metadatos del PDF...). No escribir nunca estos datos "a mano" en otros
% archivos: usar siempre estos comandos.
% =========================================================================

% --- Trabajo ---
\newcommand{\thesisTitle}{BUKIT: red social orientada a objetivos}% TODO: título definitivo
\newcommand{\thesisSubtitle}{Memoria del Trabajo de Fin de Grado}%
\newcommand{\thesisDegree}{Grado en Ingeniería Informática de Gestión y Sistemas de Información}%
\newcommand{\thesisCourse}{2026--2027}% curso académico
\newcommand{\thesisDate}{Bilbao, \today}% TODO: <Lugar, día, mes, año> de la entrega/defensa
\newcommand{\thesisYear}{2027}% año de elaboración y defensa

% --- Personas (la plantilla del centro usa "Apellido1 Apellido2, Nombre") ---
\newcommand{\thesisAuthor}{Liviu Deleanu}% nombre natural (para metadatos del PDF)
\newcommand{\thesisAuthorFormal}{Deleanu, Liviu}% TODO: confirmar formato con dos apellidos
\newcommand{\thesisDirector}{Villamañe Gironés, Mikel}%
\newcommand{\thesisCoDirector}{Apellido1 Apellido2, Nombre}% TODO: solo si hay codirector/a

% --- Universidad ---
\newcommand{\thesisUniversity}{Universidad del País Vasco / Euskal Herriko Unibertsitatea}%
\newcommand{\thesisFaculty}{Escuela de Ingeniería de Bilbao}%

% --- Palabras clave del resumen trilingüe ---
\newcommand{\thesisKeywordsEs}{red social, objetivos, monorepo, TypeScript, PostgreSQL}% TODO
\newcommand{\thesisKeywordsEu}{sare soziala, helburuak, monorepo, TypeScript, PostgreSQL}% TODO: revisar euskera
\newcommand{\thesisKeywordsEn}{social network, goals, monorepo, TypeScript, PostgreSQL}% TODO

% =========================================================================
% Comandos propios. Definir aquí cualquier cosa que se repita en el texto:
% si mañana cambia cómo se escribe, se cambia en un único sitio.
% =========================================================================

% Nombre del producto, siempre con el mismo formato.
\newcommand{\bukit}{\textsc{Bukit}\xspace}

% Términos en inglés, en cursiva de forma consistente: \eng{backend}
\newcommand{\eng}[1]{\textit{#1}}

% Notas de trabajo, visibles en rojo en el PDF. Buscar "todonote" y
% eliminarlas todas antes de entregar.
\newcommand{\todonote}[1]{\textcolor{red}{\textbf{[TODO: #1]}}}


% --- Cabecera y pie (diseño de referencia, sin logos) ---
% La caja de cabecera replica la del .tex de referencia (celda ancha de
% contexto + celda estrecha de identidad). La celda derecha llevaba el logo
% del centro: aquí va texto porque los logos UPV/EHU siguen pendientes de
% la fase de diseño (ver style.tex); cuando existan, sustituir por
% \includegraphics. 11 + 4,9 = 15,9 cm < 16 cm de texto: sin desbordes.
\usepackage{fancyhdr}
\newcommand{\riesgosHeader}{%
  \setlength{\tabcolsep}{0pt}%
  \renewcommand{\arraystretch}{1}%
  \begin{tabular}{|p{11cm}|p{4.9cm}|}
    \hline
    \parbox[c][2cm][c]{11cm}{\centering\bfseries
      {\fontsize{16}{19}\selectfont TRABAJO DE FIN DE GRADO\\[4pt]Curso \thesisCourse}%
    } &
    \parbox[c][2cm][c]{4.9cm}{\centering
      {\fontsize{14}{17}\selectfont\bfseries\bukit}\\[3pt]
      {\small\thesisFaculty}%
    }\\
    \hline
  \end{tabular}%
}
\pagestyle{fancy}
\fancyhf{}
\renewcommand{\headrulewidth}{0pt}
\fancyhead[C]{\riesgosHeader}
\fancyfoot[C]{\thepage}
% La primera página usa estilo "plain": se redefine para que también lleve
% la cabecera (mismo truco que el .tex de referencia).
\fancypagestyle{plain}{%
  \fancyhf{}%
  \renewcommand{\headrulewidth}{0pt}%
  \fancyhead[C]{\riesgosHeader}%
  \fancyfoot[C]{\thepage}%
}

% --- Símbolos y campos de formulario del diseño de referencia ---
\newcommand{\fieldline}[1]{\underline{\hspace{#1}}}
\newcommand{\bigline}{\noindent\rule{\textwidth}{0.4pt}}
\newcommand{\checkbox}{\(\square\)}
\newcommand{\checkedbox}{\(\blacksquare\)}
% Símbolo ⟳ = "si se rellena este apartado, nueva versión del plan".
% El original lo resolvía con un SVG externo (paquete svg + Inkscape);
% aquí basta el símbolo de amssymb: cero dependencias nuevas.
\newcommand{\updatesymbol}{\raisebox{-0.2ex}{\(\circlearrowright\)}}
% Resaltado amarillo del identificador de riesgo en los títulos.
\newcommand{\markfield}[1]{\colorbox{yellow!35}{\strut #1}}

% --- Enlaces y metadatos del PDF (al final, como exige hyperref) ---
% El pdftitle lo fija cada documento con \hypersetup tras cargar esto.
\usepackage[hidelinks]{hyperref}
\hypersetup{pdfauthor={\thesisAuthor}}

\hypersetup{pdftitle={Seguimiento de riesgos — 26-07-2026 (S1)}}

\begin{document}

\begin{center}
  {\LARGE\bfseries Seguimiento de riesgos — S1}\\[6pt]
\end{center}

\textbf{Fecha:} 26-07-2026 · \textbf{Sprint:} S1 (20-07 \(\to\) 26-07-2026, primera retrospectiva del proyecto) · \textbf{Responsable:} \thesisAuthor

Regla de revisión: \textbf{críticos} cada sprint; \textbf{significativos} cada 2 sprints; \textbf{moderados} cada 4 sprints; cualquier riesgo, \textbf{de inmediato}, si se observa su indicador clave.

\textbf{El símbolo \updatesymbol\ indica que, si se rellena ese apartado, hay que emitir una nueva versión completa del plan de riesgos} en \texttt{../plan-riesgos/} y registrar el cambio en la tabla de seguimiento del cap.~7 de la memoria.

Riesgos revisados: \textbf{R1–R7 por cadencia} (críticos, semanal) y \textbf{R10 por disparador observado} (su indicador clave saltó: más de tres días consecutivos sin actividad).
R8 se comenta dentro de R10 (el dato de horas de S1 se explica por la materialización de R10, no por un error estructural de capacidad; su revisión por cadencia toca en el cierre de S2).

\bigline

\section{Medidas de mitigación (prevención)}

Para cada medida de los riesgos que tocan revisión en este seguimiento: ¿se está aplicando?
Si no, justificar y decidir si se mantiene/modifica/elimina.

\begin{tabularx}{\textwidth}{|C{0.9cm}|Y|Y|C{2.2cm}|L{2.4cm}|}
\hline
\multicolumn{1}{|C{0.9cm}|}{\textbf{Id}} &
\multicolumn{1}{>{\centering\arraybackslash}X|}{\textbf{Medida de mitigación}} &
\multicolumn{1}{>{\centering\arraybackslash}X|}{\textbf{¿Aplicada?}} &
\multicolumn{1}{C{2.2cm}|}{\textbf{Justificación (del NO)}} &
\multicolumn{1}{C{2.4cm}|}{\textbf{Decisión}} \\
\hline
R1 & Toda idea nueva entra por el backlog. & \checkedbox\ Sí (la épica de vídeo EP16 entró hoy por el backlog del cap.~10, no al sprint).\newline \checkbox\ No & — & \checkedbox\ Mantener\newline \checkbox\ Modificar\newline \checkbox\ Eliminar \\
\hline
R2 & Medir velocidad real desde el primer sprint. & \checkedbox\ Sí (S1 cierra con su primera medición: 3 SP).\newline \checkbox\ No & — & \checkedbox\ Mantener\newline \checkbox\ Modificar\newline \checkbox\ Eliminar \\
\hline
R3 & Diario al cierre de cada sesión; memoria con hueco en todos los sprints. & \checkedbox\ Sí (única sesión de S1 = hoy, con diario hoy; caps.~1, 2 y 8 escritos).\newline \checkbox\ No & — & \checkedbox\ Mantener\newline \checkbox\ Modificar\newline \checkbox\ Eliminar \\
\hline
R4 & Citar al escribir; pendientes con nota visible; observar el número de páginas. & \checkedbox\ Sí (salvo la vigilancia de páginas, que desaparece: ver \updatesymbol).\newline \checkbox\ No & — & \checkbox\ Mantener\newline \checkedbox\ Modificar\newline \checkbox\ Eliminar \\
\hline
R5 & Revisar todo lo generado antes de integrarlo; nada cerrado sin comprobaciones. & \checkedbox\ Sí (verificación adversarial de caps.~1-2-8; CI en verde antes de cerrar US-01).\newline \checkbox\ No & — & \checkedbox\ Mantener\newline \checkbox\ Modificar\newline \checkbox\ Eliminar \\
\hline
R6 & Investigación previa acotada antes de tecnología desconocida. & \checkedbox\ Sí (no se tocó tecnología nueva: la CI es terreno conocido de TrackBites).\newline \checkbox\ No & — & \checkedbox\ Mantener\newline \checkbox\ Modificar\newline \checkbox\ Eliminar \\
\hline
R7 & Minimización desde el diseño (RGPD). & \checkedbox\ Sí (sin modelo de datos aún; nada que minimizar todavía).\newline \checkbox\ No & — & \checkedbox\ Mantener\newline \checkbox\ Modificar\newline \checkbox\ Eliminar \\
\hline
R10 & Margen temporal deliberado + documentación continua que permite retomar sin pérdida. & \checkedbox\ Sí (retomar hoy consistió en leer la última entrada del diario; el margen absorbió la semana).\newline \checkbox\ No & — & \checkedbox\ Mantener\newline \checkbox\ Modificar\newline \checkbox\ Eliminar \\
\hline
\end{tabularx}

\Needspace*{6\baselineskip}
\underline{\textbf{Cambios realizados en las medidas de mitigación}} \updatesymbol

\begin{tabularx}{\textwidth}{|C{1.4cm}|Y|Y|Y|}
\hline
\multicolumn{1}{|C{1.4cm}|}{\textbf{Id}} &
\multicolumn{1}{>{\centering\arraybackslash}X|}{\textbf{Medida original}} &
\multicolumn{1}{>{\centering\arraybackslash}X|}{\textbf{Cambio realizado}} &
\multicolumn{1}{>{\centering\arraybackslash}X|}{\textbf{Justificación}} \\
\hline
R4 & «El número de páginas se observa en cada compilación» (cláusula final de la mitigación). & Se elimina la cláusula; el riesgo queda centrado en la citación (plagio) y el resto de la normativa. & El autor excluyó el 26-07-2026 el apartado de extensión (40/60 pp.) de la guía de estilo aplicable a este TFG; sin límite normativo de páginas, vigilarlas ya no mitiga nada. \\
\hline
\end{tabularx}

\Needspace*{6\baselineskip}
\underline{\textbf{Nuevas medidas de mitigación a añadir}} \updatesymbol

\begin{tabularx}{\textwidth}{|C{1.4cm}|Y|Y|}
\hline
\multicolumn{1}{|C{1.4cm}|}{\textbf{Id}} &
\multicolumn{1}{>{\centering\arraybackslash}X|}{\textbf{Nueva medida}} &
\multicolumn{1}{>{\centering\arraybackslash}X|}{\textbf{Justificación}} \\
\hline
— & — & — \\
\hline
\end{tabularx}

\bigline

\section{Riesgos a revisar en este seguimiento}

% ------------------ R1 ------------------
\Needspace*{12\baselineskip}
\subsection*{Riesgo: \markfield{R1 — Crecimiento descontrolado del alcance}}

\textbf{Indicadores (clave y de resolución)}

\begin{tabularx}{\textwidth}{|Z|Z|Z|Z|Z|}
\hline
\textbf{Indicador} & \textbf{Valor esperado} & \textbf{Valor real} & \textbf{¿Cumplido?} & \textbf{Acción / observaciones} \\
\hline
Sprint gana >2 SP tras la planificación / función fuera del backlog. & 0 & 0 (EP16 entró por el backlog como opcional). & \checkbox\ Sí\quad \checkedbox\ No & Sin acción. \\
\hline
\end{tabularx}

Si el \textbf{indicador clave} se cumple: reevaluar probabilidad/impacto y severidad.
Si el \textbf{indicador de resolución} se cumple: cerrar el riesgo.

\textbf{Nueva probabilidad (1–3):} \fieldline{0.8cm} · \textbf{Nuevo impacto (1–3):} \fieldline{0.8cm} · \textbf{Nueva severidad:} \checkbox\ Crítico · \checkbox\ Significativo · \checkbox\ Moderado\ \updatesymbol

\textbf{¿Cerrar el riesgo?} \checkbox\ Sí · \checkedbox\ No\ \updatesymbol

\textbf{¿Se ha materializado?} \checkbox\ Sí · \checkedbox\ No\ \updatesymbol

\textbf{Nuevo estado} (cap.~7): \checkedbox\ Sin novedad · \checkbox\ Activo · \checkbox\ Materializado · \checkbox\ Resuelto · \checkbox\ Cerrado

\textbf{Plan de contingencia (solo si se ha materializado)}

\begin{tabularx}{\textwidth}{|Y|Y|Y|}
\hline
\multicolumn{1}{|>{\centering\arraybackslash}X|}{\textbf{Medida de contingencia}} &
\multicolumn{1}{>{\centering\arraybackslash}X|}{\textbf{¿Aplicada?}} &
\multicolumn{1}{>{\centering\arraybackslash}X|}{\textbf{Justificación (del NO)}} \\
\hline
 & \checkbox\ Sí\quad \checkbox\ No & \\
\hline
\end{tabularx}

Reflexión sobre el funcionamiento de la contingencia: no aplicable.

% ------------------ R2 ------------------
\Needspace*{12\baselineskip}
\subsection*{Riesgo: \markfield{R2 — Error sistemático de estimación del esfuerzo}}

\textbf{Indicadores (clave y de resolución)}

\begin{tabularx}{\textwidth}{|Z|Z|Z|Z|Z|}
\hline
\textbf{Indicador} & \textbf{Valor esperado} & \textbf{Valor real} & \textbf{¿Cumplido?} & \textbf{Acción / observaciones} \\
\hline
Velocidad de los 3 primeros sprints desviada >25~\% de la hipótesis. & — (necesita 3 sprints). & Primera medición: 3 SP en S1; la US-01 (3 SP) consumió \(\approx\)2,5~h (\(\approx\)0,8~h/SP frente a la hipótesis de 2,7). & \checkbox\ Sí\quad \checkedbox\ No & Un solo dato no es tendencia; seguir midiendo. \\
\hline
\end{tabularx}

Si el \textbf{indicador clave} se cumple: reevaluar probabilidad/impacto y severidad.
Si el \textbf{indicador de resolución} se cumple: cerrar el riesgo.

\textbf{Nueva probabilidad (1–3):} \fieldline{0.8cm} · \textbf{Nuevo impacto (1–3):} \fieldline{0.8cm} · \textbf{Nueva severidad:} \checkbox\ Crítico · \checkbox\ Significativo · \checkbox\ Moderado\ \updatesymbol

\textbf{¿Cerrar el riesgo?} \checkbox\ Sí · \checkedbox\ No\ \updatesymbol

\textbf{¿Se ha materializado?} \checkbox\ Sí · \checkedbox\ No\ \updatesymbol

\textbf{Nuevo estado} (cap.~7): \checkedbox\ Sin novedad · \checkbox\ Activo · \checkbox\ Materializado · \checkbox\ Resuelto · \checkbox\ Cerrado

\textbf{Plan de contingencia (solo si se ha materializado)}

\begin{tabularx}{\textwidth}{|Y|Y|Y|}
\hline
\multicolumn{1}{|>{\centering\arraybackslash}X|}{\textbf{Medida de contingencia}} &
\multicolumn{1}{>{\centering\arraybackslash}X|}{\textbf{¿Aplicada?}} &
\multicolumn{1}{>{\centering\arraybackslash}X|}{\textbf{Justificación (del NO)}} \\
\hline
 & \checkbox\ Sí\quad \checkbox\ No & \\
\hline
\end{tabularx}

Reflexión sobre el funcionamiento de la contingencia: no aplicable.

% ------------------ R3 ------------------
\Needspace*{12\baselineskip}
\subsection*{Riesgo: \markfield{R3 — Memoria relegada frente al desarrollo}}

\textbf{Indicadores (clave y de resolución)}

\begin{tabularx}{\textwidth}{|Z|Z|Z|Z|Z|}
\hline
\textbf{Indicador} & \textbf{Valor esperado} & \textbf{Valor real} & \textbf{¿Cumplido?} & \textbf{Acción / observaciones} \\
\hline
Capítulo afectado 2 sprints sin actualizar / diario >48~h tras la sesión. & 0 & 0 (diario de hoy escrito hoy; caps.~1, 2 y 8 redactados; la memoria va por delante del código). & \checkbox\ Sí\quad \checkedbox\ No & Sin acción. \\
\hline
\end{tabularx}

Si el \textbf{indicador clave} se cumple: reevaluar probabilidad/impacto y severidad.
Si el \textbf{indicador de resolución} se cumple: cerrar el riesgo.

\textbf{Nueva probabilidad (1–3):} \fieldline{0.8cm} · \textbf{Nuevo impacto (1–3):} \fieldline{0.8cm} · \textbf{Nueva severidad:} \checkbox\ Crítico · \checkbox\ Significativo · \checkbox\ Moderado\ \updatesymbol

\textbf{¿Cerrar el riesgo?} \checkbox\ Sí · \checkedbox\ No\ \updatesymbol

\textbf{¿Se ha materializado?} \checkbox\ Sí · \checkedbox\ No\ \updatesymbol

\textbf{Nuevo estado} (cap.~7): \checkedbox\ Sin novedad · \checkbox\ Activo · \checkbox\ Materializado · \checkbox\ Resuelto · \checkbox\ Cerrado

\textbf{Plan de contingencia (solo si se ha materializado)}

\begin{tabularx}{\textwidth}{|Y|Y|Y|}
\hline
\multicolumn{1}{|>{\centering\arraybackslash}X|}{\textbf{Medida de contingencia}} &
\multicolumn{1}{>{\centering\arraybackslash}X|}{\textbf{¿Aplicada?}} &
\multicolumn{1}{>{\centering\arraybackslash}X|}{\textbf{Justificación (del NO)}} \\
\hline
 & \checkbox\ Sí\quad \checkbox\ No & \\
\hline
\end{tabularx}

Reflexión sobre el funcionamiento de la contingencia: no aplicable.

% ------------------ R4 ------------------
\Needspace*{12\baselineskip}
\subsection*{Riesgo: \markfield{R4 — Incumplimiento de la normativa de la memoria}}

\textbf{Indicadores (clave y de resolución)}

\begin{tabularx}{\textwidth}{|Z|Z|Z|Z|Z|}
\hline
\textbf{Indicador} & \textbf{Valor esperado} & \textbf{Valor real} & \textbf{¿Cumplido?} & \textbf{Acción / observaciones} \\
\hline
Afirmación no original sin cita detectada en revisión. & 0 & 0 sin resolver (la verificación adversarial de los caps.~1-2 detectó 4 afirmaciones externas sin cita o sin matiz y todas se corrigieron hoy mismo). & \checkbox\ Sí\quad \checkedbox\ No & El mecanismo de verificación por capítulo funciona. \\
\hline
\end{tabularx}

Si el \textbf{indicador clave} se cumple: reevaluar probabilidad/impacto y severidad.
Si el \textbf{indicador de resolución} se cumple: cerrar el riesgo.

\textbf{Nueva probabilidad (1–3):} \fieldline{0.8cm} · \textbf{Nuevo impacto (1–3):} \fieldline{0.8cm} · \textbf{Nueva severidad:} \checkbox\ Crítico · \checkbox\ Significativo · \checkbox\ Moderado\ \updatesymbol

\textbf{¿Cerrar el riesgo?} \checkbox\ Sí · \checkedbox\ No\ \updatesymbol

\textbf{¿Se ha materializado?} \checkbox\ Sí · \checkedbox\ No\ \updatesymbol

Cambio de \textbf{plan} (no de evaluación): mitigación ajustada por la eliminación del límite de extensión (ver sección 1 \updatesymbol).
P, I y severidad no cambian.

\textbf{Nuevo estado} (cap.~7): \checkedbox\ Sin novedad (el cambio es del plan, no del estado) · \checkbox\ Activo · \checkbox\ Materializado · \checkbox\ Resuelto · \checkbox\ Cerrado

\textbf{Plan de contingencia (solo si se ha materializado)}

\begin{tabularx}{\textwidth}{|Y|Y|Y|}
\hline
\multicolumn{1}{|>{\centering\arraybackslash}X|}{\textbf{Medida de contingencia}} &
\multicolumn{1}{>{\centering\arraybackslash}X|}{\textbf{¿Aplicada?}} &
\multicolumn{1}{>{\centering\arraybackslash}X|}{\textbf{Justificación (del NO)}} \\
\hline
 & \checkbox\ Sí\quad \checkbox\ No & \\
\hline
\end{tabularx}

Reflexión sobre el funcionamiento de la contingencia: no aplicable.

% ------------------ R5 ------------------
\Needspace*{12\baselineskip}
\subsection*{Riesgo: \markfield{R5 — Dependencia acrítica de la IA generativa}}

\textbf{Indicadores (clave y de resolución)}

\begin{tabularx}{\textwidth}{|Z|Z|Z|Z|Z|}
\hline
\textbf{Indicador} & \textbf{Valor esperado} & \textbf{Valor real} & \textbf{¿Cumplido?} & \textbf{Acción / observaciones} \\
\hline
Decisión/bloque integrado que no puede explicarse sin la herramienta. & 0 & 0 (las decisiones de hoy —CI en un solo job, exclusión de e2e, retirada de \texttt{prisma generate}— están razonadas en el diario y son defendibles). & \checkbox\ Sí\quad \checkedbox\ No & Sin acción. \\
\hline
\end{tabularx}

Si el \textbf{indicador clave} se cumple: reevaluar probabilidad/impacto y severidad.
Si el \textbf{indicador de resolución} se cumple: cerrar el riesgo.

\textbf{Nueva probabilidad (1–3):} \fieldline{0.8cm} · \textbf{Nuevo impacto (1–3):} \fieldline{0.8cm} · \textbf{Nueva severidad:} \checkbox\ Crítico · \checkbox\ Significativo · \checkbox\ Moderado\ \updatesymbol

\textbf{¿Cerrar el riesgo?} \checkbox\ Sí · \checkedbox\ No\ \updatesymbol

\textbf{¿Se ha materializado?} \checkbox\ Sí · \checkedbox\ No\ \updatesymbol

\textbf{Nuevo estado} (cap.~7): \checkedbox\ Sin novedad · \checkbox\ Activo · \checkbox\ Materializado · \checkbox\ Resuelto · \checkbox\ Cerrado

\textbf{Plan de contingencia (solo si se ha materializado)}

\begin{tabularx}{\textwidth}{|Y|Y|Y|}
\hline
\multicolumn{1}{|>{\centering\arraybackslash}X|}{\textbf{Medida de contingencia}} &
\multicolumn{1}{>{\centering\arraybackslash}X|}{\textbf{¿Aplicada?}} &
\multicolumn{1}{>{\centering\arraybackslash}X|}{\textbf{Justificación (del NO)}} \\
\hline
 & \checkbox\ Sí\quad \checkbox\ No & \\
\hline
\end{tabularx}

Reflexión sobre el funcionamiento de la contingencia: no aplicable.

% ------------------ R6 ------------------
\Needspace*{12\baselineskip}
\subsection*{Riesgo: \markfield{R6 — Curva de aprendizaje de la pila tecnológica}}

\textbf{Indicadores (clave y de resolución)}

\begin{tabularx}{\textwidth}{|Z|Z|Z|Z|Z|}
\hline
\textbf{Indicador} & \textbf{Valor esperado} & \textbf{Valor real} & \textbf{¿Cumplido?} & \textbf{Acción / observaciones} \\
\hline
Historia\slash investigación al doble de su tiempo previsto. & 0 & 0 (US-01 dentro de lo previsto; sin tecnología nueva esta semana). & \checkbox\ Sí\quad \checkedbox\ No & Sin acción. \\
\hline
\end{tabularx}

Si el \textbf{indicador clave} se cumple: reevaluar probabilidad/impacto y severidad.
Si el \textbf{indicador de resolución} se cumple: cerrar el riesgo.

\textbf{Nueva probabilidad (1–3):} \fieldline{0.8cm} · \textbf{Nuevo impacto (1–3):} \fieldline{0.8cm} · \textbf{Nueva severidad:} \checkbox\ Crítico · \checkbox\ Significativo · \checkbox\ Moderado\ \updatesymbol

\textbf{¿Cerrar el riesgo?} \checkbox\ Sí · \checkedbox\ No\ \updatesymbol

\textbf{¿Se ha materializado?} \checkbox\ Sí · \checkedbox\ No\ \updatesymbol

\textbf{Nuevo estado} (cap.~7): \checkedbox\ Sin novedad · \checkbox\ Activo · \checkbox\ Materializado · \checkbox\ Resuelto · \checkbox\ Cerrado

\textbf{Plan de contingencia (solo si se ha materializado)}

\begin{tabularx}{\textwidth}{|Y|Y|Y|}
\hline
\multicolumn{1}{|>{\centering\arraybackslash}X|}{\textbf{Medida de contingencia}} &
\multicolumn{1}{>{\centering\arraybackslash}X|}{\textbf{¿Aplicada?}} &
\multicolumn{1}{>{\centering\arraybackslash}X|}{\textbf{Justificación (del NO)}} \\
\hline
 & \checkbox\ Sí\quad \checkbox\ No & \\
\hline
\end{tabularx}

Reflexión sobre el funcionamiento de la contingencia: no aplicable.

% ------------------ R7 ------------------
\Needspace*{12\baselineskip}
\subsection*{Riesgo: \markfield{R7 — Incumplimiento del RGPD en el diseño}}

\textbf{Indicadores (clave y de resolución)}

\begin{tabularx}{\textwidth}{|Z|Z|Z|Z|Z|}
\hline
\textbf{Indicador} & \textbf{Valor esperado} & \textbf{Valor real} & \textbf{¿Cumplido?} & \textbf{Acción / observaciones} \\
\hline
Dato personal sin finalidad declarada en el modelo / requisito que exija almacenar el DNI. & 0 & 0 (el modelo de datos aún no existe). & \checkbox\ Sí\quad \checkedbox\ No & Primera revisión real cuando se diseñe la entidad de usuario (US-04). \\
\hline
\end{tabularx}

Si el \textbf{indicador clave} se cumple: reevaluar probabilidad/impacto y severidad.
Si el \textbf{indicador de resolución} se cumple: cerrar el riesgo.

\textbf{Nueva probabilidad (1–3):} \fieldline{0.8cm} · \textbf{Nuevo impacto (1–3):} \fieldline{0.8cm} · \textbf{Nueva severidad:} \checkbox\ Crítico · \checkbox\ Significativo · \checkbox\ Moderado\ \updatesymbol

\textbf{¿Cerrar el riesgo?} \checkbox\ Sí · \checkedbox\ No\ \updatesymbol

\textbf{¿Se ha materializado?} \checkbox\ Sí · \checkedbox\ No\ \updatesymbol

\textbf{Nuevo estado} (cap.~7): \checkedbox\ Sin novedad · \checkbox\ Activo · \checkbox\ Materializado · \checkbox\ Resuelto · \checkbox\ Cerrado

\textbf{Plan de contingencia (solo si se ha materializado)}

\begin{tabularx}{\textwidth}{|Y|Y|Y|}
\hline
\multicolumn{1}{|>{\centering\arraybackslash}X|}{\textbf{Medida de contingencia}} &
\multicolumn{1}{>{\centering\arraybackslash}X|}{\textbf{¿Aplicada?}} &
\multicolumn{1}{>{\centering\arraybackslash}X|}{\textbf{Justificación (del NO)}} \\
\hline
 & \checkbox\ Sí\quad \checkbox\ No & \\
\hline
\end{tabularx}

Reflexión sobre el funcionamiento de la contingencia: no aplicable.

% ------------------ R10 ------------------
\Needspace*{12\baselineskip}
\subsection*{Riesgo: \markfield{R10 — Indisponibilidad prolongada del autor}}

\textbf{Indicadores (clave y de resolución)}

\begin{tabularx}{\textwidth}{|Z|Z|Z|Z|Z|}
\hline
\textbf{Indicador} & \textbf{Valor esperado} & \textbf{Valor real} & \textbf{¿Cumplido?} & \textbf{Acción / observaciones} \\
\hline
>3 días consecutivos sin poder trabajar por causa sobrevenida. & 0 & 8 días (18-07 \(\to\) 25-07) por agotamiento del trabajo de verano y asuntos personales. & \checkedbox\ Sí\quad \checkbox\ No & Revisión inmediata dentro de esta retrospectiva. \\
\hline
\end{tabularx}

Si el \textbf{indicador clave} se cumple: reevaluar probabilidad/impacto y severidad.
Si el \textbf{indicador de resolución} se cumple: cerrar el riesgo.

\textbf{Nueva probabilidad (1–3):} \underline{\ 1\ } · \textbf{Nuevo impacto (1–3):} \underline{\ 3\ } · \textbf{Nueva severidad:} \checkbox\ Crítico · \checkedbox\ Significativo · \checkbox\ Moderado\ \updatesymbol\ (sin cambio: el episodio confirma la evaluación, no la empeora)

\textbf{¿Cerrar el riesgo?} \checkbox\ Sí · \checkedbox\ No\ \updatesymbol

\textbf{¿Se ha materializado?} \checkedbox\ Sí · \checkbox\ No\ \updatesymbol

\textbf{Nuevo estado} (cap.~7): \checkbox\ Sin novedad · \checkbox\ Activo · \checkbox\ Materializado · \checkedbox\ Resuelto (se materializó y ya está superado) · \checkbox\ Cerrado

\textbf{Plan de contingencia (solo si se ha materializado)}

\begin{tabularx}{\textwidth}{|Y|Y|Y|}
\hline
\multicolumn{1}{|>{\centering\arraybackslash}X|}{\textbf{Medida de contingencia}} &
\multicolumn{1}{>{\centering\arraybackslash}X|}{\textbf{¿Aplicada?}} &
\multicolumn{1}{>{\centering\arraybackslash}X|}{\textbf{Justificación (del NO)}} \\
\hline
Pausa formal + re-planificación del calendario (+ cambio de convocatoria en el extremo). & \checkedbox\ Sí, en su versión mínima: pausa de facto y re-encaje dentro del propio S1.\newline \checkbox\ No & La re-planificación completa no fue necesaria: la historia del sprint (US-01) se completó el último día y el margen del calendario (cap.~13) absorbe el resto sin tocar el camino crítico. \\
\hline
\end{tabularx}

Reflexión sobre el funcionamiento de la contingencia:
la mitigación (margen deliberado + documentación continua) hizo su trabajo: retomar tras 8 días costó una lectura del diario, y la semana perdida cabe en el margen declarado sin re-planificar.
La contingencia prevista era proporcionalmente mayor que el episodio; se registra que su primer uso real fue la versión mínima.
Si el patrón se repitiera en S2, tocará además revisar R8 (capacidad) con dos semanas de datos.

\bigline

\section{Reflexión desde el último seguimiento}

\textbf{¿Ha ocurrido algo reseñable desde el último seguimiento?} \checkedbox\ Sí · \checkbox\ No

Es la primera retrospectiva del proyecto: S1 transcurrió sin actividad de lunes a sábado (materialización de R10) y se recuperó en la sesión de hoy, que cierra la historia del tema del sprint (US-01, CI en verde) y adelanta memoria (caps.~1, 2 y 8).
Velocidad de S1: 3 SP.

\Needspace*{6\baselineskip}
\textbf{Nuevos riesgos identificados} \updatesymbol

\begin{tabularx}{\textwidth}{|C{1.4cm}|Y|C{0.8cm}|C{0.8cm}|C{1cm}|C{2.4cm}|}
\hline
\multicolumn{1}{|C{1.4cm}|}{\textbf{Id}} &
\multicolumn{1}{>{\centering\arraybackslash}X|}{\textbf{Descripción}} &
\multicolumn{1}{C{0.8cm}|}{\textbf{P}} &
\multicolumn{1}{C{0.8cm}|}{\textbf{I}} &
\multicolumn{1}{C{1cm}|}{\textbf{Exp.}} &
\multicolumn{1}{C{2.4cm}|}{\textbf{Severidad}} \\
\hline
— & — & — & — & — & — \\
\hline
\end{tabularx}

\textbf{¿Afecta algo a los objetivos o al alcance (cap.~3 de la memoria)?} \checkbox\ Sí · \checkedbox\ No

\textbf{Justificación:} la semana perdida cabe en el margen declarado del calendario (cap.~13); el alcance no se toca (EP16 entra como opcional, fuera del compromiso).
La eliminación del límite de extensión afecta a la \emph{forma} de la memoria, no a los objetivos.

\textbf{Otros eventos relevantes:} el autor excluyó el apartado de extensión (40/60 pp.) de la guía de estilo aplicable (26-07-2026) \(\to\) cambio del plan de R4 (\updatesymbol, recogido en la v1); primera ejecución verde de la CI y primer ciclo completo tablero \(\to\) commit \(\to\) \texttt{Closes \#1}.

\end{document}
